\documentclass[11pt,a4paper]{article}

\usepackage{main-sss}
\usepackage{corrige-sss}

\renewcommand{\matiere}{Physique-Chimie}
\renewcommand{\examen}{DNB}
\renewcommand{\annee}{2026}
\renewcommand{\localisation}{Amérique du Nord}
\renewcommand{\titresujet}{Formule 1}
\renewcommand{\niveau}{Brevet}
\renewcommand{\typedocument}{Corrigé}
\renewcommand{\identifiantsujet}{26GENSCAN1}

\begin{document}

\pagination
\titre

\q[1 point]
Du point A au point B, le mouvement de la voiture est rectiligne et accéléré.

\q[1 point]
La molécule d'octane contient 8 atomes de carbone (C) et 18 atomes d'hydrogène (H).

\q[1 point]
La combustion produit du dioxyde de carbone, de formule chimique CO\textsubscript{2}.

\q[2 points]
Lors d'une transformation chimique, la masse se conserve. La masse totale des réactifs est donc égale à la masse totale des produits : $m_\text{octane} + m_\text{dioxygène} = m_\text{dioxyde de carbone} + m_\text{eau}$.

On en déduit : $m_\text{eau} = m_\text{octane} + m_\text{dioxygène} - m_\text{dioxyde de carbone}$ et donc : $m_\text{eau} = 100 + 351 - 309 = \SI{142}{kg}$.

La voiture a rejeté \SI{142}{kg} d'eau dans l'atmosphère lors de la course.

\q[1,5 points]
Par lecture graphique, pour \SI{1200}{kWh} d'énergie, le volume de dihydrogène nécessaire est de \SI{500}{L}, c'est à dire une masse d'environ \SI{36}{kg} de dihydrogène.

\q[1 point]
Remplacer l'essence par du dihydrogène présente l'avantage de réduire la masse de carburant, mais l'inconvénient de nécessiter un volume de stockage beaucoup plus important.

\q[2,5 points]
La combustion de \SI{100}{kg} d'essence produit \SI{309}{kg} de CO\textsubscript{2}, soit \SI{309000}{g}.

Pour que la fabrication du dihydrogène émette moins de CO\textsubscript{2} que la combustion de l'essence, il faut que la centrale électrique rejette moins de \SI{309000}{g} de CO\textsubscript{2} pour produire \SI{2400}{kWh}, soit une émission inférieure à : $\frac{309\,000}{2\,400} = \SI{128,75}{g/kWh}$.

Seules les centrales hydraulique, nucléaire et éolienne respectent ce critère.

\end{document}